\scanpage{655}\mspage{43}
instructif, bien sûr, si $M$ est stable par petites
$\varinjlim$, que ses objets soient accessibles etc, et
que $W$ soit \unsure{L-accessible}. (i.e. stable par
petites $\varinjlim$ filtrantes, et de plus accessible
\ldots). Le premier candidat pour $I$ auquel
je pense serait un $B_G$ ($G$ un groupe, p.ex.
$G=\pm1$), c'est en effet l'exemple aux antipodes
des ensembles ordonnés, lesquels, eux, ne
semblent pas faire \unsure{trop} de \unsure{problème|problèmes} (moyennant
l'hypothèse noethérienne \deleted{\unsure{tout}} au moins).

Les exemples suivants ne \deleted{\unsure{donnent}} sont pas
aussi convaincants que ça, hélas !

\uline{Exemple 3.10.} Soit $M$ une catégorie
quelconque, alors les couples $(\mathrm{Fl}\,M,\mathrm{Is}\,M)$
et $(\mathrm{Is}\,M,\mathrm{Fl}\,M)$ sont tous deux des
couples de Quillen
\[
 \mathrm{Fl}\,M\mathrel{\longleftarrow\mkern-9mu|\mkern-9mu\longrightarrow}\mathrm{Is}\,M
 \qquad \mathrm{Is}\,M\mathrel{\longleftarrow\mkern-9mu|\mkern-9mu\longrightarrow}\mathrm{Fl}\,M.\tag*{(3.10.1)}
\]
Il s'ensuit que si on pose
\[
 W=\mathrm{Is}\,M,\qquad \deleted{\unsure{E|F}}C=F=\mathrm{Fl}\,M,\tag*{(3.10.2)}
\]
$(W,C,F)$ est un \unsure{triple|triplet} de Quillen clos\textsuperscript{$*$}.
% slanted
{\reversemarginpar\margin{du moins si on admet que $M$ est stable par $\varinjlim$ et $\varprojlim$ \uline{finies}}}

D'ailleurs pour tout $I$ dans Cat, on a
$W(I)=\mathrm{Is}(M(I))$, donc \unsure{!}
\[
 \mathrm{Ho}_W(I)=M(I).\tag*{(3.10.3)}
\]

Donc la question de l'existence des foncteurs
$u^W_!,u^W_*$ se réduit à celle des foncteurs
$u_!,u_*$ habituels, adjoints \inserted{à g. et à dr.} de $u^*:M(I')\unsure{-|\to}M(I)$
% archive sheet number: 654

